\documentclass[12pt]{article}
\usepackage[utf8]{inputenc}
\usepackage[T1]{fontenc}
\usepackage[margin=1in]{geometry}
\usepackage{setspace}
\usepackage{xcolor}
\usepackage[colorlinks=true, urlcolor=blue!60!black, linkcolor=blue!60!black]{hyperref}
\setlength{\parindent}{2em}
\setlength{\parskip}{1.25em}
\renewcommand{\baselinestretch}{1.0}
\begin{document}
\begin{center}
{\LARGE \textit{Board of Trustees of the University of Alabama v. Garrett}, 531 U.S. 356 (2001)\par}
\vspace{0.8em}
{\large Case Brief --- Four-Page Edition\par}
\end{center}
\vspace{1em}

\section{Memory Jogger}

The Eleventh Amendment bars state employees from suing their state employers in federal court for money damages under Title I of the Americans with Disabilities Act (ADA), because Congress failed to document a pattern of unconstitutional state employment discrimination and enacted remedies not ``congruent and proportional'' under Section 5 of the Fourteenth Amendment (\href{https://www.courtlistener.com/opinion/118407/board-of-trustees-of-univ-of-ala-v-garrett/}{CourtListener}).

\section{Detailed Case Facts}

Patricia Garrett, a registered nurse, was a Director of Nursing at the University of Alabama in Birmingham Hospital. After breast-cancer treatment, she returned in July 1995 to be told she would have to give up her Director position, and she transferred to a lower-paying nurse-manager role. Milton Ash, a security officer for the Alabama Department of Youth Services, asked that his duties be modified because of chronic asthma and later sought daytime shifts because of sleep apnea. The Department denied every request.

Title I of the ADA bars covered employers, expressly including States, from discriminating against a ``qualified individual with a disability,'' requires ``reasonable accommodations'' absent ``undue hardship,'' and forbids practices with a disparate impact on disabled individuals. Congress expressly sought to abrogate state sovereign immunity. By 2001, however, \textit{Seminole Tribe} had barred abrogation under Article I and \textit{City of Boerne} had required Section 5 legislation to show ``congruence and proportionality.''

\section{Procedural History}

The Northern District of Alabama granted summary judgment for the State in both suits. The Eleventh Circuit consolidated the appeals and reversed under circuit precedent. The Supreme Court granted certiorari to resolve a circuit split, reversed as to Title I on February 21, 2001, and dismissed as improvidently granted the question concerning Title II.

\section{Judicial Votes}

Reversed, \textbf{5--4}, along the Rehnquist Court's familiar federalism alignment.
\begin{itemize}
\item \textbf{Majority:} Chief Justice Rehnquist, joined by Justices O'Connor, Scalia, Kennedy, and Thomas.
\item \textbf{Concurrence:} Justice Kennedy, joined by Justice O'Connor.
\item \textbf{Dissent:} Justice Breyer, joined by Justices Stevens, Souter, and Ginsburg.
\end{itemize}

\section{Holding}

Private federal-court suits by state employees for money damages under Title I of the ADA are barred by the Eleventh Amendment. Where the right enforced receives only rational-basis review, as disability classifications do under \textit{City of Cleburne v. Cleburne Living Center}, Congress must (1) identify a history and pattern of unconstitutional discrimination by the \textit{States themselves}, and (2) enact remedies congruent and proportional to those violations. Title I's standards still bind the States through federal suits and \textit{Ex parte Young} injunctions (footnote 9).

\section{Analysis of Opinions}

\textbf{Majority (Rehnquist).} Because Article I cannot abrogate immunity, the only question was whether Title I was ``appropriate legislation'' under Section 5. First, under \textit{Cleburne} disability is not a suspect class, so state action toward disabled persons is valid if rational; negative attitudes or fear, standing alone, ``does not a constitutional violation make.'' Second, the record showed no pattern of unconstitutional \textit{state} employment discrimination: local governments lack immunity and were excluded, and the record held only ``half a dozen examples'' plus unexamined Task Force anecdotes, mostly about public services. Third, even assuming a pattern, Title I's accommodation duty, its shifting of the ``undue hardship'' burden to the State, and its disparate-impact rule (contrary to \textit{Washington v. Davis}) far exceed what equal protection requires. Upholding abrogation would let Congress ``rewrite'' \textit{Cleburne}.

\textbf{Kennedy (concurring).} The ADA is a ``milestone on the path to a more decent, tolerant, progressive society.'' But charging a sovereign State with purposeful equal-protection violations is grave; had States systematically violated the Fourteenth Amendment, ``confirming judicial documentation'' would exist.

\textbf{Breyer (dissenting).} Congress compiled a record of ``massive, society-wide discrimination,'' and ``a legislature is not a court of law'' that must adjudicate each anecdote. Rational-basis deference is a rule of \textit{judicial} restraint, not a limit on Congress; using it to strike down a federal statute stands the principle ``on its head.'' Section 5 lets Congress use any rational means.

\section{Examples: Future Applications}

\begin{enumerate}
\item \textbf{Same Side (Immunity Bars Suit): Genetic Predisposition.} Damages suits against States for genetic-status discrimination, on a record of private employers and four unverified state anecdotes, fail: genetic status gets only rational-basis review, so abrogation fails.
\item \textbf{Same Side (Immunity Bars Suit): Self-Care Leave.} Congress lets state employees sue for damages when denied leave for their own serious medical condition, without documenting sex-stereotyped denials by States. Absent a targeted record, immunity holds.
\item \textbf{Opposite Side (Abrogation Valid): Courthouse Access.} Title II damages suits to enforce disabled citizens' fundamental right of physical access to state courts survive, as \textit{Tennessee v. Lane} (2004) later held, because a fundamental right triggers heightened scrutiny.
\item \textbf{Opposite Side (Abrogation Valid): Family-Care Leave.} A record showing state employers gave maternity leave to women but denied parallel leave to men on gender stereotypes supports abrogation, as in \textit{Nevada Department of Human Resources v. Hibbs} (2003).
\item \textbf{Fence-Sitter: Accessible University Websites.} Damages suits for inaccessible state-university course portals rest on 60 documented complaints and a narrower remedy than Title I, but the class gets only rational-basis review.
\end{enumerate}

\section{Critique}

\textbf{Progressive and structural.} Critics call \textit{Garrett} ``conservative judicial activism'' that converts an express grant of legislative power into a judicially policed record test, and strips state employees of back pay and damages, leaving only federal enforcement or \textit{Ex parte Young} injunctions.

\textbf{Textualist and doctrinal.} The Eleventh Amendment's text bars only suits by citizens of other States or foreign states, not suits by a State's own citizens like Garrett and Ash; \textit{Hans} and \textit{Seminole Tribe} rest on an atextual theory. And fusing \textit{Cleburne} with \textit{Boerne} creates a catch-22: because almost any state rule can be hypothesized as rational, Congress can almost never prove a pattern for non-suspect classes.

\section{Key Quotations}

\begin{quote}
``Thus, the result of \textit{Cleburne} is that States are not required by the Fourteenth Amendment to make special accommodations for the disabled, so long as their actions toward such individuals are rational. They could quite hardheadedly---and perhaps hardheartedly---hold to job-qualification requirements which do not make allowance for the disabled. If special accommodations for the disabled are to be required, they have to come from positive law and not through the Equal Protection Clause.''
\end{quote}
\noindent --- Chief Justice Rehnquist, Opinion of the Court, 531 U.S. at 367--368.

\begin{quote}
``Congress is the final authority as to desirable public policy, but in order to authorize private individuals to recover money damages against the States, there must be a pattern of discrimination by the States which violates the Fourteenth Amendment, and the remedy imposed by Congress must be congruent and proportional to the targeted violation. Those requirements are not met here, and to uphold the Act's application to the States would allow Congress to rewrite the Fourteenth Amendment law laid down by this Court in \textit{Cleburne}. Section 5 does not so broadly enlarge congressional authority.''
\end{quote}
\noindent --- Chief Justice Rehnquist, Opinion of the Court, 531 U.S. at 374.

\begin{quote}
``One of the undoubted achievements of statutes designed to assist those with impairments is that citizens have an incentive, flowing from a legal duty, to develop a better understanding, a more decent perspective, for accepting persons with impairments or disabilities into the larger society. The law works this way because the law can be a teacher. So I do not doubt that the Americans with Disabilities Act of 1990 will be a milestone on the path to a more decent, tolerant, progressive society.''
\end{quote}
\noindent --- Justice Kennedy, Concurring, 531 U.S. at 375.

\end{document}
